\section{Background and Heritage}
\label{sec:background}

The 2024 to 2026 wave of inter-agent protocols is not an
unprecedented invention. It is the third generation of an
agent-communication-language (ACL) lineage that began in the
mid-1990s, was largely abandoned during the service-oriented
architecture era, and has been recapitulated, often without
attribution, atop the modern JSON-RPC 2.0 and OpenAPI
substrates. Reading the contemporary specifications against their
heritage reveals why some design choices recurred verbatim
(performative-style message types, capability descriptors,
matchmaker-style discovery) and why others (full ontology
commitments, layered semantics in the FIPA-ACL sense) have not
returned despite long-standing arguments in their favor. This
section locates each axis of the rubric in the prior art whose
absence the LLM-era protocols are now re-discovering.

\subsection{FIPA-ACL and KQML: what worked, what didn't}
\label{sec:background:fipa-kqml}

KQML emerged from the DARPA Knowledge Sharing Effort in the
early 1990s and became the first widely-used agent communication
language~\cite{finin1997kqml}. Its central insight was the
separation of communicative \emph{performatives} (tell, ask,
subscribe, achieve) from message content. A message was a
nested expression. The outermost form indicated the
illocutionary force of the communication: whether the sender
was asserting a fact, requesting an action, or committing to a
future obligation. The inner content remained
language-agnostic. The Foundation for Intelligent Physical
Agents (FIPA) standardized a refined successor, FIPA-ACL,
in 2002~\cite{fipa_acl}, anchoring the performative vocabulary
in Searle's speech-act theory~\cite{searle1969speechacts} and
adding interaction protocols (request, contract-net, query)
that prescribed legal performative sequences.

What worked in FIPA-ACL and KQML, and what carries through to the
2024 to 2026 wave, is the recognition that an inter-agent message
carries more than payload: it carries an \emph{illocutionary
contract} between sender and receiver. A message tagged as a
request commits the sender to interpret a refusal differently
than the same payload tagged as an inform. Modern A2A messages,
ACP request envelopes, and MCP tool-call responses all carry an
implicit performative tag (typically encoded in a method name or
event type field) that performs the same function. The
contract-net interaction protocol has task announcement, bidding,
award, and completion phases~\cite{smith1980contractnet}. It
recurs in ACNBP's ten-step capability
negotiation and binding flow~\cite{spec_acnbp}, in Coral's team
contracts~\cite{spec_coral}, and (more implicitly) in A2A's
task lifecycle states~\cite{spec_a2a}. FIPA adopted it as a
named interaction protocol; the 2024 to 2026 protocols
re-derive the phase structure without citing it.

What did not work, and is not returning at the protocol layer,
was the full ontology commitment. FIPA-ACL specified that the
content language (typically FIPA-SL) and the ontology of
predicates and terms be declared in every message, with the
expectation that participating agents share a common ontology
or invoke a translation service. In practice, ontology
maintenance proved unsustainable: cross-organizational
deployments stalled on ontology mismatches that no available
matchmaker could resolve, and the practical FIPA deployments of
the late 1990s and early 2000s eventually shifted to schemaful
but human-mediated message designs. The contemporary protocols
have re-discovered the gap. Section~\ref{sec:semantic} treats
this in depth: A2A's free-text \texttt{Skill.description}
field~\cite{spec_a2a}, MCP's natural-language tool descriptions,
and ANP's JSON-LD vocabulary~\cite{spec_anp} occupy three
different points on the spectrum between unconstrained text
(easy authoring, brittle interoperability) and full
ontology commitment (heavy authoring, machine-tractable
matching). None of them re-adopts FIPA-SL.

\subsection{Service-oriented heritage: SOAP/WSDL \(\rightarrow\) REST \(\rightarrow\) OpenAPI}
\label{sec:background:soa}

Between the decline of FIPA in the mid-2000s and the rise of
LLM-era agents, inter-system communication consolidated around
the service-oriented stack. SOAP, with its XML envelope and
WSDL service descriptors, attempted to be the universal RPC
substrate but suffered from the same maintenance cost that
sank FIPA: descriptor proliferation, ontology drift, tooling
brittleness. REST emerged as the lighter-weight
alternative~\cite{fielding2002rest}, trading machine-readable
contract precision for human-readable URL conventions and
HTTP verb semantics. OpenAPI (formerly
Swagger) restored a machine-readable contract layer atop REST
without imposing SOAP's envelope overhead~\cite{spec_openapi}.

The 2024 to 2026 inter-agent protocols inherit two structural
choices from this lineage. First, all of them are HTTP-based
except where DID resolution and on-chain settlement
mandate alternative substrates. A2A runs over HTTPS+JSON-RPC,
MCP over stdio or HTTPS+SSE, ACP over REST+multipart, and ANP
over HTTPS+JSON-LD~\cite{spec_a2a, spec_mcp, spec_acp,
spec_anp}. Second, the practical contract format is JSON
Schema or its equivalents: MCP tool descriptors carry
JSON-Schema input definitions, A2A AgentCards carry typed
Skill objects, ACP carries multipart MIME types for richer
payloads. None of the protocols seriously revisits XML.

JSON-RPC 2.0~\cite{spec_jsonrpc} in particular has emerged as
the lingua franca of the tool-binding and peer-messaging
layers. MCP adopted it
directly; A2A specifies JSON-RPC 2.0 as the primary transport
with optional gRPC and REST bindings. The choice is not
arbitrary: JSON-RPC carries enough envelope structure (id,
method, params, result, error) to support correlated
request-response semantics without the schema-language
complexity of SOAP, and its compatibility with HTTP and
WebSocket transports allows the same message format to flow
through synchronous calls, streaming subscriptions, and
batched bulk operations.

\subsection{The LLM disruption: why old ACLs are returning}
\label{sec:background:llm-disruption}

The structural disruption introduced by LLM-based agents is
that the sender of a message and the receiver of a message can,
for the first time, share a flexible interlingua: natural
language. That interlingua absorbs schema mismatches that would have
crashed a FIPA-era agent. An LLM-based receiver can interpret a
loosely-typed request, identify ambiguities, and either resolve
them by inference or ask a clarifying question. The strict
contract precision that FIPA-SL demanded is partially obviated:
two LLM agents with reasonable defaults and natural-language
fluency can interoperate over a thin envelope without explicit
ontology alignment. This is why the contemporary protocols are
willing to ship with free-text tool and skill descriptions
that would have been considered fatal under FIPA: the receiver
is presumed to handle the schema impedance through reasoning.

The flip side is that this flexibility re-introduces every
attack surface that FIPA-era specifications either treated
formally or did not have to face. Prompt injection in tool
inputs~\cite{zhan2024injecagent} has no FIPA analogue because
FIPA-era content was parsed deterministically. Capability
escalation through ambiguous descriptions, supply-chain attacks
through registry poisoning, and replay attacks against
loosely-versioned tools were either out of scope or were
handled by infrastructure layers (X.509 PKI for service
authentication, WS-Security for SOAP envelopes) that the
contemporary protocols have not uniformly re-adopted.
Section~\ref{sec:security} catalogues the resulting threat
surface in detail.

The LLM disruption thus simultaneously reduces the need for
ontological strictness at the protocol layer and increases the
need for security primitives that the protocol layer did not
historically have to specify. The 2024 to 2026 protocols have
inherited the loosely-typed message style of LLM tool calling
and are now retroactively bolting on the identity, delegation,
audit, and threat-modeling primitives that FIPA-era systems
either took for granted (because they ran on closed corporate
networks) or did not have to face (because their interpreters
were deterministic). The rubric of \S\ref{sec:rubric} is
constructed to render these retrofits comparable across
protocols.

\subsection{Terminology and scope}
\label{sec:background:terms}

We adopt the following terminology uniformly throughout the
survey, departing from the more permissive usage in the prior
surveys where context demands precision.

\paragraph{Agent.} A computational entity that can produce or
consume messages on an inter-agent protocol and that holds an
identity distinguishable to other entities on that protocol.
The agent need not be an LLM; the rubric scores the protocol,
not the implementation behind a protocol endpoint.

\paragraph{Protocol.} A specification governing the syntax,
semantics, and sequencing of messages exchanged between agents.
We include only protocols with published written specifications;
implementation-defined behaviors of specific agent frameworks
(LangGraph, AutoGen, ADK) are not protocols in this sense.

\paragraph{Message.} A single unit of protocol-level
communication, with envelope metadata distinguishable from
payload content. We use ``message'' for both synchronous
request/response units and asynchronous event-stream chunks
where the envelope semantics are equivalent.

\paragraph{Task.} A first-class lifecycle-bearing object that
groups a sequence of related messages with shared state and a
terminal status. A2A's \texttt{Task}~\cite{spec_a2a}, ACNBP's
bound execution unit~\cite{spec_acnbp}, and Coral's persistent
thread~\cite{spec_coral} are tasks in this sense; MCP's
historically stateless request/response pattern is not, though
the emerging MCP Tasks extension begins to add task semantics.

\paragraph{Capability.} A named, typed, optionally
parameterized operation that an agent declares it can perform.
MCP tools, A2A skills, ACNBP capability descriptors, and Coral
tools (via MCP integration) are capabilities in this sense.
We distinguish capabilities (what the protocol declares the
agent can do) from \emph{intents} (what a calling agent
indicates it wants done), the latter recurring as a primitive
in ANP and LOKA but not in tool-binding-style protocols.

\paragraph{Principal.} The human, organizational, or
machine identity on whose behalf an agent acts. Principal
delegation is the cryptographic or
policy-enforced statement that an agent is authorized to act
on behalf of a named principal. It becomes a first-class concern
in the commerce-settlement layer (AP2, Visa TAP, ERC-8004) and
in the security analysis of Sections~\ref{sec:security}
and~\ref{sec:governance}.
